\chapter{Results}



\section{Effect of \acs*{ArBIB} on the microstructure of hydrated cementitious specimens}

The \ac{ArBIB} process enables to produce very clean surfaces of specimens.% as shown in Section~\ref{sec:bib_pol_fixed}.
However, it causes some artifacts, which affect the evaluation of the resulting images.
Therefore, this chapter takes a closer look at the artifacts caused by \ac{ArBIB}.


As indicated in Figure~\ref{fig:bib_scheme}, the ion beam can induce heat in the surface of the specimen during milling and polishing.
To determine the influence of the ion beam during the regular milling process on the actual structure of the hydration products, it hs to be compared with a milling process while cooling the specimen.
The ion beam device used, allows to cool the specimen below {\color{red} \SI{-140}{\celsius}} using liquid nitrogen during the static milling process (Section~\ref{sec:bib_pol_fixed}).

The comparison was made using hydrated alite, hydrated \ac{OPC} and \ac{LC3}.

\dots

A well known artifact